% =====================================================================
% Revisão de literatura (Ensaio 1 / dissertação PPGEco-UFES)
% Versão pré-triagem, 2026-10-12. Arquivo autocontido para o Overleaf:
% a bibliografia está embutida abaixo (filecontents) e é gerada na
% compilação. Compilar com pdfLaTeX + BibTeX (padrão do Overleaf).
% Fonte no repositório: paper/overleaf/revisao-literatura.tex
% =====================================================================
\begin{filecontents*}[overwrite]{refs.bib}
@article{dias2023down,
  author  = {Dias, Mateus and Rocha, Rudi and Soares, Rodrigo R},
  title   = {Down the River: Glyphosate Use in Agriculture and Birth Outcomes of Surrounding Populations},
  journal = {Review of Economic Studies},
  year    = {2023},
  volume  = {90},
  number  = {6},
  pages   = {2943--2981},
  doi     = {10.1093/restud/rdad011}
}

@article{reynier2025glyphosate,
  author  = {Reynier, Emmett and Rubin, Edward},
  title   = {Glyphosate exposure and {GM} seed rollout unequally reduced perinatal health},
  journal = {Proceedings of the National Academy of Sciences},
  year    = {2025},
  volume  = {122},
  number  = {3},
  doi     = {10.1073/pnas.2413013121}
}

@article{camacho2017health,
  author  = {Camacho, Adriana and Mejía, Daniel},
  title   = {The health consequences of aerial spraying illicit crops: The case of Colombia},
  journal = {Journal of Health Economics},
  year    = {2017},
  volume  = {54},
  pages   = {147--160},
  doi     = {10.1016/j.jhealeco.2017.04.005}
}

@article{larsen2017agricultural,
  author  = {Larsen, Ashley E. and Gaines, Steven D. and Deschênes, Olivier},
  title   = {Agricultural pesticide use and adverse birth outcomes in the San Joaquin Valley of California},
  journal = {Nature Communications},
  year    = {2017},
  volume  = {8},
  number  = {1},
  doi     = {10.1038/s41467-017-00349-2}
}

@article{frank2024ecosystem,
  author  = {Frank, Eyal G.},
  title   = {The economic impacts of ecosystem disruptions: Costs from substituting biological pest control},
  journal = {Science},
  year    = {2024},
  volume  = {385},
  number  = {6713},
  doi     = {10.1126/science.adg0344}
}

@techreport{callaway2024continuous,
  author  = {Callaway, Brantly and Goodman-Bacon, Andrew and Sant'Anna, Pedro H.},
  title   = {Difference-in-differences with a Continuous Treatment},
  institution = {National Bureau of Economic Research},
  type    = {NBER Working Paper},
  number  = {32117},
  year    = {2024},
  doi     = {10.3386/w32117}
}

@article{weitzman1974,
  author  = {Weitzman, Martin L.},
  title   = {Prices vs. Quantities},
  journal = {The Review of Economic Studies},
  year    = {1974},
  volume  = {41},
  number  = {4},
  pages   = {477},
  doi     = {10.2307/2296698}
}

@article{helfand1995,
  author  = {Helfand, Gloria E. and House, Brett W.},
  title   = {Regulating Nonpoint Source Pollution Under Heterogeneous Conditions},
  journal = {American Journal of Agricultural Economics},
  year    = {1995},
  volume  = {77},
  number  = {4},
  pages   = {1024--1032},
  doi     = {10.2307/1243825}
}

@article{lichtenberg1988,
  author  = {Lichtenberg, Erik and Parker, Douglas D. and Zilberman, David},
  title   = {Marginal Analysis of Welfare Costs of Environmental Policies: The Case of Pesticide Regulation},
  journal = {American Journal of Agricultural Economics},
  year    = {1988},
  volume  = {70},
  number  = {4},
  pages   = {867--874},
  doi     = {10.2307/1241928}
}

@article{skevas2012,
  author  = {Skevas, Theodoros and Stefanou, Spiro E. and Lansink, Alfons Oude},
  title   = {Can economic incentives encourage actual reductions in pesticide use and environmental spillovers?},
  journal = {Agricultural Economics},
  year    = {2012},
  volume  = {43},
  number  = {3},
  pages   = {267--276},
  doi     = {10.1111/j.1574-0862.2012.00581.x}
}

@article{bocker2020glyphosate,
  author  = {Böcker, Thomas and Britz, Wolfgang and Möhring, Niklas and Finger, Robert},
  title   = {An economic and environmental assessment of a glyphosate ban for the example of maize production},
  journal = {European Review of Agricultural Economics},
  year    = {2019},
  note    = {Advance access. Volume, numero e paginas da versao final nao constam do deposito Crossref},
  doi     = {10.1093/erae/jby050}
}

@article{finger2024europe,
  author  = {Finger, Robert},
  title   = {Europe's ambitious pesticide policy and its impact on agriculture and food systems},
  journal = {Agricultural Economics},
  year    = {2024},
  volume  = {55},
  number  = {2},
  pages   = {265--269},
  doi     = {10.1111/agec.12817}
}

@article{arkhangelsky2021sdid,
  author  = {Arkhangelsky, Dmitry and Athey, Susan and Hirshberg, David A. and Imbens, Guido W. and Wager, Stefan},
  title   = {Synthetic {Difference-in-Differences}},
  journal = {American Economic Review},
  year    = {2021},
  volume  = {111},
  number  = {12},
  pages   = {4088--4118},
  doi     = {10.1257/aer.20190159}
}

@article{conley2011,
  author  = {Conley, Timothy G. and Taber, Christopher R.},
  title   = {Inference with {``Difference} in Differences'' with a Small Number of Policy Changes},
  journal = {Review of Economics and Statistics},
  year    = {2011},
  volume  = {93},
  number  = {1},
  pages   = {113--125},
  doi     = {10.1162/rest_a_00049}
}

@article{canay2021wild,
  author  = {Canay, Ivan A. and Santos, Andres and Shaikh, Azeem M.},
  title   = {The Wild Bootstrap with a {``Small''} Number of {``Large''} Clusters},
  journal = {The Review of Economics and Statistics},
  year    = {2021},
  volume  = {103},
  number  = {2},
  pages   = {346--363},
  doi     = {10.1162/rest_a_00887}
}

@article{shirangi2010living,
  author  = {Shirangi, Adeleh and Nieuwenhuijsen, Mark and Vienneau, Danielle and Holman, C.D'Arcy J.},
  title   = {Living near agricultural pesticide applications and the risk of adverse reproductive outcomes: a review of the literature},
  journal = {Paediatric and Perinatal Epidemiology},
  year    = {2010},
  volume  = {25},
  number  = {2},
  pages   = {172--191},
  doi     = {10.1111/j.1365-3016.2010.01165.x}
}

@article{calzada2023bananas,
  author  = {Calzada, Joan and Gisbert, Meritxell and Moscoso, Bernard},
  title   = {The Hidden Cost of Bananas: The Effects of Pesticides on Newborns’ Health},
  journal = {Journal of the Association of Environmental and Resource Economists},
  year    = {2023},
  volume  = {10},
  number  = {6},
  pages   = {1623--1663},
  doi     = {10.1086/725349}
}

@article{neginegi2025misclassified,
  author  = {Negi, Akanksha and Negi, Digvijay S.},
  title   = {Difference‐in‐Differences With a Misclassified Treatment},
  journal = {Journal of Applied Econometrics},
  year    = {2025},
  volume  = {40},
  number  = {4},
  pages   = {411--423},
  doi     = {10.1002/jae.3116}
}

@article{ling2018prenatal,
  author  = {Ling, Chenxiao and Liew, Zeyan and von Ehrenstein, Ondine S. and Heck, Julia E. and Park, Andrew S. and Cui, Xin and Cockburn, Myles and Wu, Jun and Ritz, Beate},
  title   = {Prenatal Exposure to Ambient Pesticides and Preterm Birth and Term Low Birthweight in Agricultural Regions of {California}},
  journal = {Toxics},
  year    = {2018},
  volume  = {6},
  number  = {3},
  pages   = {41},
  doi     = {10.3390/toxics6030041}
}

@article{vanwendel2024isa,
  author  = {van Wendel de Joode, Berna and Peñaloza-Castañeda, Jorge and Mora, Ana María and Corrales-Vargas, Andrea and Eskenazi, Brenda and Hoppin, Jane A. and Lindh, Christian H.},
  title   = {Pesticide exposure, birth size, and gestational age in the {ISA} birth cohort, {Costa Rica}},
  journal = {Environmental Epidemiology},
  year    = {2024},
  volume  = {8},
  number  = {2},
  pages   = {e290},
  doi     = {10.1097/EE9.0000000000000290}
}

@article{kadhel2014chlordecone,
  author  = {Kadhel, Philippe and Monfort, Christine and Costet, Nathalie and Rouget, Florence and Thomé, Jean-Pierre and Multigner, Luc and Cordier, Sylvaine},
  title   = {Chlordecone exposure, length of gestation, and risk of preterm birth},
  journal = {American Journal of Epidemiology},
  year    = {2014},
  volume  = {179},
  number  = {5},
  pages   = {536--544},
  doi     = {10.1093/aje/kwt313}
}

@article{herve2016chlordecone,
  author  = {Hervé, David and Costet, Nathalie and Kadhel, Philippe and Rouget, Florence and Monfort, Christine and Thomé, Jean-Pierre and Multigner, Luc and Cordier, Sylvaine},
  title   = {Prenatal exposure to chlordecone, gestational weight gain, and birth weight in a {Guadeloupean} birth cohort},
  journal = {Environmental Research},
  year    = {2016},
  volume  = {151},
  pages   = {436--444},
  doi     = {10.1016/j.envres.2016.08.004}
}

@article{rouget2020chlordecone,
  author  = {Rouget, Florence and Kadhel, Philippe and Monfort, Christine and Viel, Jean-François and Thomé, Jean-Pierre and Cordier, Sylvaine and Multigner, Luc},
  title   = {Chlordecone exposure and risk of congenital anomalies: the {Timoun} Mother-Child Cohort Study in {Guadeloupe} ({French West Indies})},
  journal = {Environmental Science and Pollution Research},
  year    = {2020},
  volume  = {27},
  number  = {33},
  pages   = {40992--40998},
  doi     = {10.1007/s11356-019-06031-y}
}

@article{brainerd2014seasonal,
  author  = {Brainerd, Elizabeth and Menon, Nidhiya},
  title   = {Seasonal effects of water quality: The hidden costs of the {Green Revolution} to infant and child health in {India}},
  journal = {Journal of Development Economics},
  year    = {2014},
  volume  = {107},
  pages   = {49--64},
  doi     = {10.1016/j.jdeveco.2013.11.004}
}

@misc{mejia2014enforcement,
  author       = {Mejía, Daniel and Restrepo, Pascual and Rozo, Sandra},
  title        = {On the Effects of Enforcement on Illegal Markets: Evidence from a Quasi-Experiment in {Colombia}},
  year         = {2014},
  howpublished = {SSRN Working Paper 2481408},
  doi          = {10.2139/ssrn.2481408}
}

@article{page2021prisma,
  author  = {Page, Matthew J. and McKenzie, Joanne E. and Bossuyt, Patrick M. and others},
  title   = {The {PRISMA} 2020 statement: an updated guideline for reporting systematic reviews},
  journal = {BMJ},
  year    = {2021},
  pages   = {n71},
  doi     = {10.1136/bmj.n71}
}

@article{mcgowan2016press,
  author  = {McGowan, Jessie and Sampson, Margaret and Salzwedel, Douglas M. and Cogo, Elise and Foerster, Vicki and Lefebvre, Carol},
  title   = {{PRESS} Peer Review of Electronic Search Strategies: 2015 Guideline Statement},
  journal = {Journal of Clinical Epidemiology},
  year    = {2016},
  volume  = {75},
  pages   = {40--46},
  doi     = {10.1016/j.jclinepi.2016.01.021}
}

\end{filecontents*}
\documentclass[11pt,a4paper]{article}
\usepackage[T1]{fontenc}
\usepackage[utf8]{inputenc}
\usepackage[brazil]{babel}
\usepackage{lmodern}
\usepackage{microtype}
\usepackage{amsmath}
\usepackage{booktabs}
\usepackage{xcolor}
\usepackage[margin=2.3cm]{geometry}
\usepackage[colorlinks=true,linkcolor=black,citecolor=blue!50!black,urlcolor=blue!50!black]{hyperref}
\usepackage[round,authoryear]{natbib}
\setlength{\bibsep}{2pt}
\setlength{\parskip}{3pt}

\title{\vspace{-1.5cm}Proibição da pulverização aérea no Ceará:\\ revisão de literatura}
\author{Felipe Carvalho Souza Santos\\\small PPGEco/UFES}
\date{\small Versão pré-triagem, outubro de 2026}

\begin{document}
\maketitle
\vspace{-0.6cm}
\section{Revisão de literatura}
\label{sec:revisao}

Esta seção organiza a literatura em torno de uma pergunta só: \emph{o que se
sabe sobre o dano perinatal de agrotóxicos agrícolas que permita avaliar a
proibição de um método de aplicação?} A pergunta importa porque a Lei Estadual
nº~16.820/2019 proíbe a pulverização \emph{aérea}, e não uma molécula, enquanto
quase toda a evidência mede moléculas. Quatro eixos ordenam o que segue: a
substância (o que o agrotóxico faz ao feto), o método (o que a via aérea
acrescenta), o instrumento (o que o regulador precisa saber) e o estimador (o
que um choque simultâneo com dose contínua permite identificar). A seção
termina com o que a literatura não responde, e com o que este trabalho pode e
não pode responder.

\subsection{Como a literatura foi levantada}

O levantamento segue um protocolo de revisão sistemática fixado antes da busca
\citep{page2021prisma}. Ele tem uma pergunta principal em formato PECOS:
exposição \emph{ambiental} na gestação a agrotóxico de uso agrícola e desfechos
ao nascer, com o método de aplicação como moderador. Tem também duas revisões
estruturadas complementares, uma de método e uma de instrumento. As estratégias
de busca foram calibradas contra estudos conhecidos e submetidas a uma revisão
PRESS independente \citep{mcgowan2016press}, que exigiu correções nos quatro
blocos. A Tabela~\ref{tab:busca} resume o resultado.

\begin{table}[h]
\centering\small
\caption{Estratégias de busca congeladas (2026-10-12)}
\label{tab:busca}
\begin{tabular}{@{}llrl@{}}
\toprule
Bloco & Bases & Registros brutos & Recall de itens conhecidos \\
\midrule
SQ1, internacional & OpenAlex, PubMed & 3.385 + 717 & 13/13 \\
SQ1, América Latina & OpenAlex, SciELO, LILACS & $\approx$2.300 & 21/21 (indexados) \\
SQ2, método & OpenAlex, arXiv & $\approx$4.100 & 34/38 \\
SQ3, instrumento & OpenAlex, Wiley, web & $\approx$5.000 & 18/18 (no ajuste) \\
\bottomrule
\end{tabular}
\par\smallskip
\begin{minipage}{0.92\linewidth}\footnotesize
\emph{Nota:} recall medido nos itens usados no desenvolvimento da estratégia, e
não em validação independente. A triagem ainda não foi feita. Os estudos
citados abaixo vêm do corpus já conferido do projeto e dos estudos que a
calibração encontrou, cada um conferido na fonte primária (PubMed ou Crossref).
\end{minipage}
\end{table}

A calibração já ensinou duas coisas sobre a literatura, antes mesmo da triagem.
A primeira é que os dois estudos econômicos mais próximos deste caso
\citep{calzada2023bananas,dias2023down} não estão indexados no PubMed: uma
revisão feita só em bases biomédicas não os encontraria. A segunda é que um
filtro de contexto ``agrícola e espacial'', que parecia inofensivo, eliminava
as coortes de clordecona das Antilhas Francesas. São justamente os estudos de
um inseticida de bananal. A literatura mais próxima do caso cearense não se
descreve com o vocabulário que a busca esperava.

\subsection{A substância: dano concentrado na cauda, e desigual}

O resultado mais importante da literatura para este trabalho é de forma, não de
nível. \citet{larsen2017agricultural}, com mais de meio milhão de nascimentos no
Vale de San Joaquin, não encontram efeito na média: a exposição residencial eleva
os desfechos adversos em 5 a 9\% \emph{apenas} acima do percentil 95.
\citet{reynier2025glyphosate}, explorando a adoção de sementes geneticamente
modificadas com um instrumento de aptidão agroclimática, acham redução do peso e
da duração da gestação, com efeito doze vezes maior no decil inferior de peso
esperado. \citet{dias2023down} chegam ao mesmo tipo de dano pela água: o
glifosato a montante eleva a mortalidade infantil a jusante.
\citet{frank2024ecosystem} chegam a ele pelo avesso: a morte de morcegos
predadores leva o agricultor a usar 31,1\% mais inseticida, e a mortalidade
infantil sobe 7,9\%.

A epidemiologia de exposição individual confirma o dano, mas mostra que
\emph{qual} desfecho responde depende da molécula. Na Califórnia, a exposição
residencial num raio de 2~km eleva o risco de prematuridade em 3 a 7\%, e quase
nunca o de baixo peso a termo \citep{ling2018prenatal}. Na coorte ISA, da Costa
Rica, um aumento de dez vezes no metabólito do clorpirifós se associa a 130~g a
menos de peso entre nascidos a termo, sem efeito sobre a duração da gestação
\citep{vanwendel2024isa}. Em Guadalupe, a clordecona, o inseticida usado nos
bananais até 1993, encurta a gestação em 0,27 semana e eleva o risco de
prematuridade em 60\% para cada aumento de dez vezes na concentração plasmática \citep{kadhel2014chlordecone}. A mesma coorte,
porém, não encontra associação com o peso ao nascer no agregado
\citep{herve2016chlordecone}, nem com malformações \citep{rouget2020chlordecone}.

Três regularidades saem daí. O dano está na cauda da exposição, não na média. A
margem que responde, peso ou duração da gestação, muda com a molécula. E há
resultados nulos sólidos, inclusive na população mais exposta a um agrotóxico de
bananal. Uma revisão que relatasse só os efeitos encontrados descreveria outra
literatura.

\subsection{O método: a via aérea}

A evidência sobre o \emph{método} vem de dois contextos andinos. Na Colômbia, a
aspersão aérea de herbicida sobre cultivos ilícitos foi política de Estado, e
\citet{camacho2017health} usam as campanhas por município e os registros
individuais de saúde para estimar suas consequências sanitárias. No Equador,
\citet{calzada2023bananas} estudam a fumigação aérea dos bananais e cruzam o
endereço da mãe com o perímetro das plantações e com a estação de fumigação
intensa, a chuvosa, quando a sigatoka-negra ataca. O resumo publicado reporta
déficit de 80 a 150~g no peso ao nascer sob exposição alta. É o análogo mais
próximo do caso cearense nas três dimensões que importam: cultura, método e
classe química, já que a procimidona, onipresente nas amostras da Chapada do
Apodi, é fungicida de bananal.

Um terceiro estudo mostra o instrumento do Ensaio~2 em operação.
\citet{mejia2014enforcement} exploram a suspensão da aspersão numa faixa de
10~km junto à fronteira com o Equador, imposta em 2006 por atrito diplomático.
Por descontinuidade na faixa, estimam que cada hectare pulverizado reduz o
cultivo de coca em apenas 0,02 a 0,065 hectare. Essa faixa é, na prática, uma
\emph{zona-tampão}, mas foi avaliada pelo efeito sobre o alvo da pulverização, e
não sobre a saúde de quem mora ao lado. Por fim, \citet{brainerd2014seasonal}
mostram como identificar dano sazonal: cruzam o mês da concepção com a estação
em que fertilizantes chegam à água e encontram efeito sobre a mortalidade
infantil. O químico é outro, mas a lógica de calendário é a mesma de que este trabalho dispõe se a medida do
método vier a existir.

Nenhum desses estudos avalia a \emph{retirada} de um método. Na Colômbia, a via
aérea foi imposta; no Equador, ela segue em uso. No Ceará, ela foi proibida para
todos de uma vez. Isso muda tanto o sinal esperado quanto a estrutura de quem
escolhe se expor.

\subsection{O instrumento e o estimador}

A comparação entre preço e quantidade sob incerteza depende das inclinações
relativas do dano marginal e do custo marginal \citep{weitzman1974}. A economia
dos agrotóxicos herdou essa moldura: mede o custo de bem-estar da regulação
\citep{lichtenberg1988}, trata o problema como poluição difusa sob
heterogeneidade \citep{helfand1995}, avalia taxas \citep{skevas2012} e simula
\emph{ex ante} a proibição de uma molécula \citep{bocker2020glyphosate,
finger2024europe}. Esses trabalhos supõem uma função de dano, mas nenhum estima
a inclinação dela para um \emph{método} de aplicação. A calibração da busca
apontou ainda dois ramos que faltavam ao menu original. O primeiro é a
responsabilização \emph{ex post} pela deriva, um quarto instrumento ao lado de
proibição, taxa e zona-tampão. O segundo é a elasticidade-preço da demanda por
agrotóxico, o parâmetro de custo que uma taxa exige.

Do lado econométrico, três resultados delimitam o que um choque estadual
simultâneo com dose contínua pode entregar. \citet{callaway2024continuous}
separam o efeito no nível de dose recebido, $\text{ATT}(d|d)$, identificado sob
tendências paralelas, da curva causal, que exige uma hipótese mais forte e
exclui a seleção pelo ganho. A literatura de tratamento mal classificado cobre o
erro por falso negativo \citep{neginegi2025misclassified}. Aqui o erro é o
oposto: municípios com dose positiva que não tinham avião. Para esse caso, o
resultado relevante é o de Denteh e Kédagni,\footnote{\emph{Working paper},
arXiv:2207.11890, ainda sem versão publicada; o argumento que dele depende
repousa, portanto, num texto não publicado.} segundo o qual, sob
tendências paralelas e sem falso negativo, o estimando é o efeito verdadeiro
multiplicado pelo valor preditivo positivo da classificação: o erro atenua, mas
não inverte o sinal.
E com poucos controles, a inferência precisa vir de métodos desenhados para
isso \citep{conley2011,canay2021wild,arkhangelsky2021sdid}.

\subsection{Onde este trabalho entra, e onde não entra}

A lacuna tem dois tempos, e é preciso escrevê-los separadamente para não
prometer o que o corpo não entrega. O primeiro é o da literatura: não há
avaliação \emph{ex post} da proibição de um método de aplicação, e não há
estimativa da inclinação do dano para a via aérea numa cultura não transgênica
tratada com fungicida. O segundo é o deste trabalho. Ele oferece a primeira
avaliação \emph{ex post} de uma proibição estadual da pulverização aérea, com
dose contínua. O que ela entrega é o nível, $\text{ATT}(d|d)$, e um limite
superior informativo, condicionado ao poder do desenho. A curva dose-resposta,
que a comparação de Weitzman exigiria, não é interpretável nos dados (desvio
registrado na pré-especificação em 22/09/2026). O Ensaio~2, por isso, trabalha
com limites, não com uma inclinação pontual. É menos do que a literatura de
instrumento pede, e mais do que ela tem hoje.

{\small
\bibliographystyle{plainnat}
\bibliography{refs}
}
\end{document}
